% Conclusion and future work (P8.7)
\section{Conclusion and future work}\label{sec:conclusion}

This project built a land suitability analysis for urban development in Bengaluru South
entirely in Python, from open land-cover, terrain and OpenStreetMap data, and tested every
method against the growth that actually happened. It answers the three questions it set out
with:

\begin{enumerate}
  \item \textbf{Land classes.} Clustering 100\,m cells on land-cover shares and terrain gives
        seven interpretable clusters, mapped to built-up (348\,km\textsuperscript{2} in 2023),
        forest (29\,km\textsuperscript{2}, including the steep scrub forest of Bannerghatta) and
        usable land (138\,km\textsuperscript{2}). The map agrees with the ESRI classes on
        96.6\,\% of land cells.
  \item \textbf{Building similarity.} For any building, the usable locations whose surroundings
        look most like its own can be found and explained feature by feature. Averaged over
        reference buildings, this similarity predicts growth from 2018 to 2022/23 with
        AUC 0.773, without ever seeing a growth label.
  \item \textbf{Suitability.} Combining the similarity with a random forest trained on recent
        change gives the best score (AUC 0.797; the top 10\,\% of cells hold 3.45 times their
        share of growth), clearly ahead of a classic AHP overlay (0.636) and ahead of the
        distance to built-up land (0.763). The final map grades usable land into classes
        S1--S3 and lists the 20 best candidate sites.
\end{enumerate}

The main lesson is about validation: strict rules against leakage (no own-cell land cover, data
frozen at the baseline date, out-of-fold scores) turn near-perfect accuracies into honest ones,
and simple baselines show how much of the signal is just ``next to existing built-up land''.

\paragraph{Future work.}
\begin{itemize}
  \item Add drivers of growth away from the built-up edge: planned layouts, industrial estates,
        land prices, distance to the city centre and to employment.
  \item Re-run the analysis with newer years as they become available, moving the baseline
        forward to test whether the models generalise in time.
  \item Use building footprints from other sources (e.g. Open Buildings) to type buildings, so
        that similarity can be computed per building type.
  \item Test finer grids with a growth label defined independently of the cell size.
\end{itemize}
